\section{A common-zero signature for the remaining positive excess}
\label{sec:booker-rankin}

The signed-discrepancy conclusion in MA-LD-001 can be combined
with symmetric-square holomorphy on the whole critical line.
This gives a condition on common zeros of specified functions,
beyond the inequality at their central values.

Fix the finite set $S$ used in MA-S5-006 and put
\[
 Q_{\rm sym}(s)=L^S(s,\Pi,\operatorname{Sym}^2\otimes\nu^{-1}),
 \qquad D(s)=L_{V\eta}(s)L_X(s),\qquad s_t=\tfrac12+it.
 \tag{MA-LD.12}
\]
For real $t$ define the finite meromorphic orders
\[
 \begin{aligned}
 \beta(t)&=\operatorname{ord}_{s_t}B,&
 b(t)&=\operatorname{ord}_{s_t}C,&
 n(t)&=\operatorname{ord}_{s_t}A,\\
 d(t)&=\operatorname{ord}_{s_t}L_{V\eta},&
 e(t)&=\operatorname{ord}_{s_t}L_X,&
 q(t)&=\operatorname{ord}_{s_t}Q_{\rm sym},\\
 \kappa(t)&=\operatorname{ord}_{s_t}L_W-n(t)
             =\beta(t)-b(t)-n(t).
 \end{aligned}
 \tag{MA-LD.13}
\]
All of $\beta,b,n,d,e,q$ are nonnegative. Here the functions
$B,C$ are nonprincipal Hecke functions, $A,L_{V\eta},L_X$
are the known entire automorphic functions already used in
MA-S5-001/003, and the assertion for $Q_{\rm sym}$ is the
holomorphy part of Takeda's theorem. None of these assertions
requires a zero-location hypothesis. At $t=0$,
$\kappa(0)=k$.

\paragraph{MA-LD-004 (two unavoidable common-zero families).}
At every real $t$ one has
\[
 \begin{gathered}
 \boxed{q(t)=d(t)+e(t)-\kappa(t)},\\
 -(b(t)+n(t))\le\kappa(t)\le d(t)+e(t),
 \end{gathered}
 \tag{MA-LD.14}
\]
and therefore
\[
                 \beta(t)+q(t)=b(t)+n(t)+d(t)+e(t).
 \tag{MA-LD.15}
\]
Assume now $(\mathcal R_B)$ and $k>0$. Then both of the following
sets of heights are infinite and unbounded in absolute value:
\[
 \{t\in\mathbb R:B(s_t)=D(s_t)=0\},\qquad
 \{t\in\mathbb R:Q_{\rm sym}(s_t)=A(s_t)C(s_t)=0\}.
 \tag{MA-LD.16}
\]
This assertion is about literal coincident zeros at the same
height, with their actual orders. Define their overlap counts
with multiplicity by
\[
 \begin{aligned}
 N_{B,D}(T)&=\sum_{|t|\le T}\min\{\beta(t),d(t)+e(t)\},\\
 N_{Q,AC}(T)&=\sum_{|t|\le T}\min\{q(t),b(t)+n(t)\}.
 \end{aligned}
 \tag{MA-LD.17}
\]
The sums are over the discrete union of the relevant zero
heights. Then
\[
       N_{B,D}(T)+N_{Q,AC}(T)\ge N_R(T)\ne o(T).
 \tag{MA-LD.18}
\]
Consequently, finiteness of either common-zero set in
(MA-LD.16), or an $o(T)$ bound for the sum in (MA-LD.18),
would imply $k=0$ under $(\mathcal R_B)$.

\begin{proof}
First extend the local-factor check of MA-S5-003 from the
central point to the entire line $\operatorname{Re}s=1/2$.
Every finite correction factor there has roots that are
roots of unity in the variable $p^{-s}$, so its zeros and
poles have real part zero. The archimedean corrections are
ratios of products of $\Gamma_{\mathbb R}(s)$ and
$\Gamma_{\mathbb R}(s+1)$, because the relevant parameters
have finite image. Their poles are at nonpositive real
integers. Thus these factors, as well as conductor
exponentials, are holomorphic and nonvanishing at every
$s_t$. Deleting the finitely many additional factors in
the chosen $S$ has the same property.

Write $z(t)=\operatorname{ord}_{s_t}\zeta_{\mathbb Q}$ and
$r_4(t)=\operatorname{ord}_{s_t}
\Lambda(s,\Pi\times\Pi^\vee)$.
The meromorphic identity (MA-S5.11) and the character
decomposition (MA-S5.8) therefore give
\[
 \begin{split}
 r_4(t)
 &=z(t)+d(t)+\operatorname{ord}_{s_t}L_W+e(t)
                                       -2\kappa(t)\\
 &=z(t)+n(t)+d(t)+e(t)-\kappa(t).
 \end{split}
\]
The symmetric/exterior-square identity used in MA-S5-006,
at this same point, gives
\[
                         r_4(t)=z(t)+n(t)+q(t).
\]
The orders $z(t)$ need not vanish away from the centre,
but cancel exactly between these two formulas. This proves
the identity in (MA-LD.14). Takeda's Theorem~7.1 states
holomorphy of the indicated incomplete symmetric-square
function away from the possible points $0$ and $1$;
in particular $q(t)\ge0$ for every real $t$. This proves
the upper bound. The lower bound follows from
$\beta(t)=b(t)+n(t)+\kappa(t)\ge0$, and (MA-LD.15) is
the same pair of identities combined.

Under the positive-excess hypotheses, MA-LD-001 proves
infinitely many negative and positive orders for the
completed comparison quotient, all on the critical line.
Its order at $s_t$ is $-\kappa(t)$. Hence there are
infinitely many heights with $\kappa(t)>0$, and infinitely
many with $\kappa(t)<0$. At a height of the first kind,
\[
              \beta(t)\ge\kappa(t)>0,
              \qquad d(t)+e(t)\ge\kappa(t)>0.
\]
At a height of the second kind,
\[
              q(t)\ge-\kappa(t)>0,
              \qquad b(t)+n(t)\ge-\kappa(t)>0.
\]
These prove (MA-LD.16). The functions involved are nonzero
meromorphic functions, so each bounded segment contains
only finitely many zeros; the two sets are therefore
unbounded. The same inequalities give, at every height,
\[
 \min\{\beta(t),d(t)+e(t)\}
     +\min\{q(t),b(t)+n(t)\}\ge|\kappa(t)|.
\]
By critical-line support, $N_R(T)=\sum_{|t|\le T}|\kappa(t)|$.
Summation and (MA-LD.5) prove (MA-LD.18). The final sufficient
criteria follow by contradiction.
\end{proof}

This deduction gives no smaller numerical upper bound than
$k\le d(0)+e(0)$ at the centre. Its new requirement is an
infinite set of coincidences away from the centre whose
combined multiplicity count is not $o(T)$. No theorem excluding that
pattern is asserted. The zero-location input $M$ alone
does not bound these coincidences or their multiplicities.
In particular the argument does not apply a low-degree
converse theorem to the symmetric-square function, whose
degree is ten, and it does not assume that it is an
automorphic standard function on $\mathrm{GL}_{10}$.

\paragraph{Primary input and check scope.}
The new use of Takeda is his holomorphy on the whole line,
not merely at the central point. The exact primary source
read is Theorem~7.1 and its complete terminal proof,
\texttt{arXiv:1506.04791v1}, pp.~49--50. Its permitted $S$
is exactly the archimedean places, places dividing $2$,
and the ramified places of the representation and twist.
The finite-order twist and unitary central character
were checked in MA-S5-006. The proof invokes Takeda's
Theorem~6.2 and earlier integral representation; these
upstream representation-theoretic arguments remain
published inputs rather than newly reproved results.

Takeda, S. (2015). On a certain metaplectic Eisenstein
series and the twisted symmetric square $L$-function.
\emph{Mathematische Zeitschrift, 281}, 103--157.
\url{https://doi.org/10.1007/s00209-015-1476-x}.
\url{https://arxiv.org/abs/1506.04791v1}.
